\subsection{Stock forecast benchmarks}\label{sec:stock_validation_methods}

The option forecast errors motivated a separate evaluation of the stock generator. We first checked return scaling and price reconstruction and repeated the 2026 stock forecasts after conditioning the initial state on the available return history. This diagnostic retained the fitted parameters and tracked the remaining duration of a jump episode during state inference. We then compared four stock forecast methods: the saved 2014--2024 JumpHMM, an unchanged-price benchmark, a Gaussian model with adaptive volatility, and the same volatility model with a simple directional regression. The original JumpHMM retained its stationary initialization and independent-pilot growth shift. The adaptive model updated daily variance with an exponentially weighted moving average (EWMA) of squared log returns. The regression estimated a daily drift from the stock's most recent one, five, and twenty valid daily returns and SPY's most recent five valid daily returns. Cumulative-return features were normalized by origin volatility and the square root of their window length. A ridge penalty shrank the coefficients toward zero, and the regression had no intercept. Let $u_t=\log(S_t/S_{t-1})$ denote an observed daily return and $a_t$ its estimated daily variance at the forecast origin. The adaptive and directional models used the equations:
\begin{equation}\label{eq:stock_benchmark}
a_t=\lambda a_{t-1}+(1-\lambda)u_t^2,
\qquad
\log\frac{S_{t+h}}{S_t}=h\mu_t+\sqrt{a_t}\sum_{j=1}^{h}Z_j,
\quad Z_j\sim\mathcal N(0,1).
\end{equation}
The adaptive model set $\mu_t=0$, while the directional model used the regression estimate. Drift and variance remained fixed within each forecast horizon, so neither model received future stock or market returns. The unchanged-price benchmark supplied only a point forecast. It was not assigned an uncertainty interval for comparison with the stochastic models. These stock tests did not replace the paths used in the original option evaluation or the paired IV ablation.

We selected the new models' settings using expanding annual validation blocks from 2019 through 2024, with earlier observations used for each regression fit and feature scale. One EWMA decay and one ridge penalty were chosen across both tickers by weighting each ticker, validation year, and horizon equally. Candidate decays were $0.80$, $0.90$, $0.94$, $0.97$, and $0.99$. Ridge penalties ranged from $0.01$ to $1000$ and included the zero-drift limit. Decay selection minimized Gaussian log loss for cumulative log returns, while penalty selection minimized squared cumulative-return error divided by forecast variance. After selection, we fitted the regression using 2014--2024 observations and froze its coefficients and scaling. The additional 2025 stock dataset supplied 249 one-session and 245 five-session origins per ticker. The August--September 2026 test retained the original 21 one-session and 17 five-session stock outcomes per ticker. All methods used identical starting prices and endpoints. Origin features and variance updated only with observations then available; missing closes were not filled or converted into multi-session returns treated as one-day observations. Each stochastic model used 10{,}000 paths per origin and three fixed simulation seeds, with shared Gaussian innovations between the adaptive and directional candidates. We reported median-based MAE, mean-based RMSE, CRPS, and central 90\% interval coverage and width. Medians minimize expected absolute loss, while means minimize expected squared loss~\cite{hyndman2021}. Scores were averaged across dates within each simulation replicate and then across replicates. The model classes were proposed after inspecting the original 2026 misses, but neither the 2025 nor 2026 outcomes selected their settings. We therefore reported both periods separately and treated the 2026 follow-up as exploratory.
